% Part: proof-theory
% Chapter: proof-search
% Section: search-algorithm

\documentclass[../../../include/open-logic-section]{subfiles}

\begin{document}

\olfileid[id]{pt}{ps}{sal}

\olsection{Algoritma Pencarian untuk \Log{G3c}}

Kita mendeskripsikan sebuah algoritma pencarian untuk \Log{G3c}.
Algoritma ini bukan satu-satunya algoritma yang mungkin, bahkan bukan
yang paling efisien. Namun, algoritma ini benar: jika suatu sekuen
$\Gamma \Sequent \Delta$ memiliki !!a{proof}, algoritma ini pada akhirnya
akan berhenti dan menemukan !!a{proof}. Algoritma ini juga menjamin keberhasilan:
jika ia berhenti tanpa menemukan !!a{proof}, atau tidak pernah berhenti,
maka sejak awal $\Gamma \Sequent \Delta$ memang tidak sahih.

Salah satu ciri pokok yang diperlukan algoritma ini ialah jaminan bahwa
setiap !!{formula} dalam $\Gamma \Sequent \Delta$, atau dalam sekuen yang
dihasilkan selama pencarian bukti, ``direduksi,'' yakni diperlakukan
sebagai formula utama dari suatu aturan inferensi yang diterapkan secara
mundur, dan bahwa hal ini dilakukan sesering yang diperlukan. Ciri ini
kita sebut \emph{keadilan}.

Algoritma bekerja secara bertahap. Keluaran setiap tahap adalah sebuah
pohon sekuen, dan keluaran tahap berikutnya diperoleh dengan mereduksi
!!a{formula} dalam setiap sekuen paling atas yang belum merupakan aksioma.

Pada tahap~$0$ kita mulai dengan $\Gamma \Sequent \Delta$. Untuk menjamin
keadilan, kita memberikan bilangan sebagai indeks kepada !!{formula}s
dalam $\Gamma$ dan $\Delta$ sedemikian rupa sehingga !!{formula}s yang
berbeda memperoleh bilangan yang berbeda. Misalnya, bayangkan kita
menghitung !!{formula}s dari kiri ke kanan. Bilangan tersebut kita tulis
sebagai superskrip. Sebagai contoh, jika kita mencari !!a{proof} atas
$!A, !B \Sequent !C$, keluaran tahap~$0$ ialah $!A^1, !B^2 \Sequent
!C^3$. Bilangan terbesar~$n$ yang muncul sebagai superskrip dalam sekuen
berindeks semacam itu disebut \emph{indeks maksimalnya} ($3$ dalam contoh
tersebut).

Jika sekuen memuat kuantor, kita juga perlu mengetahui suku mana yang
digunakan ketika mereduksi $\lexists[x][!A(x)]$ di kanan atau
$\lforall[x][!A(x)]$ di kiri. Kita hanya memerlukan suku yang dibangun
dari !!{constant}s $\Obj c_0$, $\Obj c_1$, $\Obj c_2$, \dots, dengan
menggunakan sembarang simbol fungsi yang muncul dalam~$\Gamma$ dan
$\Delta$. Kita mengandaikan bahwa himpunan semua suku ini diberikan dalam
suatu enumerasi tetap, sehingga kita dapat berbicara, misalnya, tentang
``!!{constant} pertama yang tidak muncul dalam $\Gamma \Sequent
\Delta$.'' Setiap sekuen berindeks dalam pohon yang dihasilkan algoritma
memiliki suatu himpunan !!{constant}s yang terkait dengannya. Himpunan
!!{constant}s yang diberikan kepada sekuen pada tahap~$0$ ialah himpunan
semua !!{constant}s yang muncul di dalamnya (jika ada), dan
$\{\Obj c_0\}$ jika sekuen itu tidak memuat !!{constant}.

Andaikan algoritma telah menghasilkan sebuah pohon sekuen berindeks,
beserta himpunan !!{constant}s yang terkait, pada tahap~$n$. Pada
tahap~$n+1$, kita melakukan hal berikut untuk setiap sekuen paling atas.

Jika ada !!a{formula}~$!A$ yang muncul baik di sisi kiri maupun kanan
sekuen, atau sisi kirinya memuat~$\lfalse$, kita tidak melakukan apa pun:
sekuen-sekuen demikian memiliki !!{proof}s dalam~\Log{G3c}, dan untuk
sekuen paling atas tersebut tidak ada lagi yang perlu dilakukan.

Jika sekuen tidak memuat !!{formula} nonatomik, hentikan algoritma dengan
kegagalan. Tidak ada bukti atas $\Gamma \Sequent \Delta$.

Jika tidak, pilih !!{formula} nonatomik~$!A$ dengan indeks terkecil~$i$.
Sekuen paling atas yang dimaksud berbentuk $!A^i, \Pi \Sequent \Lambda$
atau $\Pi \Sequent \Lambda, !A^i$. Misalkan $k$ adalah indeks maksimal
sekuen itu, dan $C$ himpunan !!{constant}s yang diberikan kepadanya.

Kita ``mereduksi'' $!A^i$, yakni menghasilkan sekuen-sekuen paling
atas yang baru menurut aturan inferensi yang ditentukan oleh
!!{main operator} dari~$!A$ serta menurut apakah $!A^i$ muncul di kiri
atau di kanan.

\begin{enumerate}
  \item $!A \ident !B \land !C$ dan muncul di kiri: Tambahkan satu sekuen
  paling atas baru di atas $(!B \land !C)^i, \Pi \Sequent \Lambda$:
  \[
  \Axiom$!B^{k+1}, !C^{k+2}, \Pi \fCenter \Lambda$
  \RightLabel{\LeftR{\land}}
  \UnaryInf$(!B \land !C)^i, \Pi \fCenter \Lambda$
  \DisplayProof
  \]
  Himpunan !!{constant}s yang diberikan kepada sekuen baru ialah~$C$.
  \item $!A \ident !B \land !C$ dan muncul di kanan: Tambahkan dua sekuen
  paling atas baru di atas $\Pi \Sequent \Lambda, (!B \land !C)^i$:
  \[
  \Axiom$\Pi \fCenter \Lambda, !B^{k+1}$
  \Axiom$\Pi \fCenter \Lambda, !C^{k+1}$
  \RightLabel{\RightR{\land}}
  \BinaryInf$\Pi \fCenter \Lambda, (!B \land !C)^i$
  \DisplayProof
  \]
  Himpunan !!{constant}s yang diberikan kepada sekuen-sekuen baru
  ialah~$C$.

  \item Kasus $!A \ident \lnot !B$, $!A \ident !B \lor !C$, dan
  $!A \ident !B \lif !C$ serupa dan diserahkan sebagai latihan.

  \item $!A \ident \lexists[x][!B(x)]$ dan muncul di kiri: Tambahkan satu
  sekuen paling atas baru di atas $\lexists[x][!B(x)]^i, \Pi \Sequent
  \Lambda$:
  \[
  \Axiom$!B(c)^{k+1}, \Pi \fCenter \Lambda$
  \RightLabel{\LeftR{\lexists}}
  \UnaryInf$\lexists[x][!B(x)]^i, \Pi \fCenter \Lambda$
  \DisplayProof
  \]
  dengan $c$ sebagai !!{constant} pertama yang tidak berada dalam~$C$.
  Himpunan !!{constant}s yang diberikan kepada sekuen baru ialah
  $C \cup \{c\}$.
  \item $!A \ident \lexists[x][!B(x)]$ dan muncul di kanan: Tambahkan satu
  sekuen paling atas baru di atas $\Pi \Sequent \Lambda,
  \lexists[x][!B(x)]^i$:
  \[
  \Axiom$\Pi \fCenter \Lambda, \lexists[x][!B(x)]^{k+1},
    !B(t)^{k+2}$
  \RightLabel{\RightR{\lexists}}
  \UnaryInf$\Pi \fCenter \Lambda, \lexists[x][!B(x)]^i$
  \DisplayProof
  \]
  dengan $t$ sebagai suku pertama yang hanya memuat !!{constant}s
  dalam~$C$ dan belum digunakan dalam reduksi
  $\lexists[x][!B(x)]$ di kanan di bawah sekuen ini (atau $\Obj c_0$
  jika semua suku tersebut telah digunakan). Himpunan !!{constant}s
  yang diberikan kepada sekuen baru ialah~$C$.
  \item Kasus $!A \ident \lforall[x][!B(x)]$ serupa dan diserahkan
  sebagai latihan.
\end{enumerate}

Jika pada tahap~$n+1$ kita tidak menambahkan sekuen baru kepada sekuen
paling atas mana pun, algoritma berhenti dengan berhasil. Dalam hal ini
kita telah menemukan !!a{proof} atas~$\Gamma \Sequent \Delta$, yang
diperoleh dari pohon pada tahap~$n+1$ dengan menghapus semua indeks.
Semua sekuen paling atas berbentuk $!A, \Pi \Sequent \Lambda, !A$ atau
$\lfalse, \Pi \Sequent \Lambda$. Semua inferensinya merupakan inferensi
yang benar dalam~\Log{G3c}. Hal ini jelas bagi inferensi proposisional.
Perhatikan bahwa pada setiap langkah, himpunan~$C$ dari !!{constant}s
yang diberikan kepada suatu sekuen memuat semua !!{constant}s yang
muncul dalam sekuen itu.
Karena itu, dalam reduksi yang menggunakan \LeftR{\lexists} dan
\RightR{\lforall}, syarat variabel eigen dipenuhi: variabel eigen~$c$
adalah !!a{constant} yang tidak berada dalam himpunan~$C$ milik
kesimpulan, sehingga $c$ juga tidak muncul dalam kesimpulan. Kita
memperoleh:

\begin{prop}
Algoritma pencarian bukti untuk \Log{G3c} adalah benar: jika algoritma
berhenti dengan berhasil, sekuen awal~$\Gamma \Sequent \Delta$ memiliki
!!a{proof}.
\end{prop}

\end{document}
